\documentclass[12pt,a4paper]{iit-project-draft-report}
\usepackage[spanish]{babel}
\usepackage[utf8]{inputenc}
\usepackage[T1]{fontenc}
\usepackage{tabularx}
\usepackage{graphicx}
\usepackage{booktabs}
\usepackage{amsmath,amssymb}
\usepackage{siunitx}
\usepackage[acronym]{glossaries}
\usepackage[nottoc]{tocbibind}
\usepackage{cite}
\usepackage{hyperref}
\hypersetup{hidelinks}

\makeglossaries

%%%%%%%%%%%%%%%%%%%%%%%%%%%%%%%%%%%%%%%%%%%%%%%%%%%%%%%%%
% Acrónimos
%%%%%%%%%%%%%%%%%%%%%%%%%%%%%%%%%%%%%%%%%%%%%%%%%%%%%%%%%
\newacronym{CA}{CA}{Corriente Alterna}
\newacronym{CC}{CC}{Corriente Continua}
\newacronym{OPF}{OPF}{flujo óptimo de potencia}
\newacronym{SOC}{SOC}{cono de segundo orden}
\newacronym{SOCP}{SOCP}{programación cónica de segundo orden}
\newacronym{SDP}{SDP}{programación semidefinida}
\newacronym{HVDC}{HVDC}{transmisión en corriente continua de alta tensión}
\newacronym{VSC}{VSC}{convertidor fuente de tensión}
\newacronym{GFL}{GFL}{Grid-Following}
\newacronym{FUBM}{FUBM}{Flexible Universal Branch Model}
\newacronym{CFBM}{CFBM}{Convex Flexible Branch Model}
\newacronym{NTC}{NTC}{capacidad neta de transferencia}
\newacronym{TSO}{TSO}{operador del sistema de transmisión}
\newacronym{ENTSOE}{ENTSO-E}{European Network of Transmission System Operators for Electricity}
\newacronym{REE}{REE}{Red Eléctrica de España}
\newacronym{RTE}{RTE}{Réseau de Transport d'Électricité}
\newacronym{Terna}{Terna}{operador italiano del sistema de transmisión}
\newacronym{MATPOWER}{MATPOWER}{paquete de simulación de sistemas eléctricos en MATLAB}
\newacronym{VFlexP}{VFlexP}{herramienta del IIT para sistemas híbridos AC/DC}

%%%%%%% COMANDOS PARA LA TABLA DE LA PORTADA
\newcommand{\projectname}{}
\newcommand{\project}[1]{\renewcommand{\projectname}{#1}}
\newcommand{\reviewernames}{}
\newcommand{\revisedby}[1]{\renewcommand{\reviewernames}{#1}}
\newcommand{\datedetailsname}{}
\newcommand{\datedetails}[1]{\renewcommand{\datedetailsname}{#1}}
\newcommand{\statename}{}
\newcommand{\state}[1]{\renewcommand{\statename}{#1}}
\newcommand{\referencenametable}{}
\newcommand{\referencetable}[1]{\renewcommand{\referencenametable}{#1}}
\newcommand{\ficheroname}{}
\newcommand{\fichero}[1]{\renewcommand{\ficheroname}{#1}}

\makeatletter

\project{CONECCMAR}
\title{Flujo óptimo de potencia mediante relajación convexa SOC para sistemas híbridos AC/DC. Aplicación al enlace HVDC España--Italia}
    \let\Title\@title
\author{Mohammad Rajabdorri}
    \let\Author\@author
\revisedby{}
\datedetails{21-05-2026}
    \let\Date\@date
\state{Borrador}
\referencetable{Sin ref. no.}
\fichero{SOCOPF\_doc.zip}

\makeatother

\begin{document}

\maketitle

\noindent \hrulefill

\tableofcontents\thispagestyle{fancy}

\noindent \hrulefill

\glsresetall

\newpage
\section{Introducción}

El proyecto CONECCMAR contempla un enlace \gls{HVDC} submarino de \SI{2}{GW}
entre España e Italia. El análisis estacionario de este tipo de sistemas
híbridos AC/DC requiere herramientas de flujo de cargas y de \gls{OPF} que
representen correctamente el modelo de los convertidores \gls{VSC} y la red
de continua. La herramienta \gls{VFlexP}, construida sobre \gls{MATPOWER}
con la extensión \gls{FUBM} \cite{AlvarezBustos2021}, resuelve el problema
mediante un esquema iterativo Newton--Raphson. Al tratarse de una
formulación no convexa no garantiza el óptimo global y, en sistemas mal
condicionados o con varios convertidores, puede converger a soluciones
locales o no converger.

Las relajaciones convexas del problema de \gls{OPF} han alcanzado un grado
de madurez notable en la última década
\cite{MolzahnHiskens2019,Low2014a,Low2014b}. Las dos familias principales
son la relajación \gls{SDP}, teóricamente la más ajustada pero costosa de
resolver, y la relajación \gls{SOCP} (en su forma de cono rotado), menos
ajustada pero significativamente más eficiente y soportada por solvers
comerciales como Gurobi. Trabajos recientes \cite{KocukDeySun2016} han
demostrado que, con refuerzos adecuados, la relajación \gls{SOCP} puede
ser tan ajustada como la \gls{SDP} en una amplia variedad de casos. Para
redes radiales se ha probado además que ambas relajaciones son
equivalentes.

El \gls{CFBM} propuesto en \cite{HigueraGutierrez2025} extiende la
relajación \gls{SOCP} clásica para incorporar los convertidores \gls{VSC}
en una formulación unificada con la red de alterna. Cada convertidor se
representa mediante una rama de transformador no ideal entre el nudo de
continua y un nudo virtual, seguida de una rama-$\pi$ entre el nudo
virtual y el nudo físico de alterna, junto con una fuente reactiva
auxiliar acotada en el nudo virtual. Esta descomposición permite mantener
el problema entero dentro de la formulación \gls{SOCP}.

El objetivo de este documento es presentar la implementación de un
\gls{OPF} basado en la relajación \gls{CFBM}-\gls{SOC} para el caso de
estudio de CONECCMAR, validarlo frente al solver Newton--Raphson de
referencia y discutir su precisión, sus restricciones activas y su uso
potencial como solver convexo independiente o como punto inicial para
métodos iterativos no convexos.

\section{Formulación del problema}
\label{sec:formulation}

Sea $\mathcal{B}$ el conjunto de nudos, $\mathcal{L}_{AC}$ y
$\mathcal{L}_{DC}$ los conjuntos de líneas de \gls{CA} y \gls{CC}
respectivamente, y $\mathcal{C}$ el conjunto de convertidores. Se introduce
el cambio de variables clásico de la relajación \gls{SOCP}: para cada nudo
$i$, $W_{ii} = |V_i|^2 \ge 0$; y para cada par $(i,j)$ conectado por una
rama, $W_{ij} = V_i V_j^{*}$, descompuesto en sus partes real e imaginaria
$W^{re}_{ij}$ y $W^{im}_{ij}$.

\subsection{Restricciones de la red}

Las ecuaciones de balance de potencia activa y reactiva en cada nudo, los
límites de tensión y las restricciones térmicas de las ramas resultan
lineales en las variables $W$. La única no linealidad procede del
acoplamiento entre $W_{ii}$, $W_{jj}$ y $W_{ij}$, que se sustituye por la
restricción de cono rotado
\begin{equation}
4 \left( (W^{re}_{ij})^2 + (W^{im}_{ij})^2 \right) +
\left( W_{ii} - W_{jj} \right)^2 \le \left( W_{ii} + W_{jj} \right)^2,
\label{eq:soc}
\end{equation}
que es la relajación cónica de la identidad
$|V_i V_j^{*}|^2 = W_{ii} W_{jj}$. La relajación es ajustada cuando
\eqref{eq:soc} se satisface con igualdad en el óptimo.

\subsection{Modelo CFBM del convertidor VSC}

Cada convertidor $k \in \mathcal{C}$ se representa, siguiendo
\cite{HigueraGutierrez2025}, mediante tres elementos:
\begin{enumerate}
  \item Una rama de transformador no ideal entre el nudo primario de
  continua $p$ y un nudo virtual $m$. En la formulación \gls{SOCP} se
  impone $W^{im}_{pm}=0$, $W^{re}_{pm}\ge 0$ y la cota
  $\tau_{min}^{2}\,W_{mm} \le W_{pp} \le \tau_{max}^{2}\,W_{mm}$
  para limitar el módulo del tap del convertidor.
  \item Una rama-$\pi$ convencional entre el nudo virtual $m$ y el nudo
  físico de alterna $a$, con parámetros $(r,x,b)$ del convertidor.
  \item Una fuente reactiva auxiliar $Q_z$ en el nudo virtual, acotada por
  \begin{equation}
  B_{eq}^{min}\,W_{mm} \le Q_z \le B_{eq}^{max}\,W_{mm}.
  \label{eq:qz}
  \end{equation}
\end{enumerate}
El paso de potencia activa se conserva
($P^{NI}_{p \to m} = -P^{NI}_{m \to p}$) y la reactiva por el lado de
continua es nula ($Q^{NI}_{p \to m} = 0$). Las pérdidas por conmutación
del convertidor se modelan mediante una conductancia $G_{sw}$ en el nudo
primario.

\subsection{Convertidores en regulación de tensión}

En este estudio los convertidores se modelan como estaciones \gls{GFL}
que regulan la tensión del nudo de \gls{CA} a \SI{1}{pu}. Esto se impone
añadiendo la restricción $W_{aa} = 1$ para cada convertidor. La potencia
activa transmitida queda libre, sujeta únicamente a los límites térmicos
del cable de continua y de la propia estación convertidora.

\subsection{Función objetivo}

Se utiliza un objetivo lineal proporcional al despacho activo de los
generadores,
\begin{equation}
\min \quad \sum_{g \in \mathcal{G}} c_{1,g}\,P_g,
\label{eq:obj}
\end{equation}
donde $c_{1,g}$ es el coste marginal del generador $g$. La elección de
estos coeficientes se discute en la sección~\ref{sec:data}.

\section{Caso de estudio}
\label{sec:case}

El caso de estudio es el sistema de prueba de nueve nudos con dos
convertidores \gls{CC}/\gls{CA} descrito en \cite{TomasMartin2026}, que
agrupa la generación y la carga de España, Francia e Italia en sendos
nudos equivalentes interconectados por dos líneas de \gls{CA} y un enlace
\gls{HVDC} bipolar. La base del sistema es $S_b = \SI{2000}{MVA}$.

\section{Datos y supuestos}
\label{sec:data}

\subsection{Capacidad instalada de generación}

Como $P_{max}$ de cada nudo-país se utiliza la capacidad despachable
instalada en 2024 publicada por los respectivos \gls{TSO} (\gls{REE},
\gls{RTE} y \gls{Terna}). El límite reactivo se fija en $\pm 30\%$ de
$P_{max}$, una envolvente PQ típica para generación agregada a escala
nacional (Tabla~\ref{tab:gencap}).

\begin{table}[h]
\centering
\caption{Capacidad instalada utilizada como $P_{max}$ y envolvente reactiva.}
\label{tab:gencap}
\begin{tabular}{lccc}
\toprule
Nudo / país & $P_{max}$ [MW] & $Q_{max}$ [MVAr] & $Q_{min}$ [MVAr] \\
\midrule
1 -- España & 120\,000 & $+36\,000$ & $-36\,000$ \\
7 -- Francia & 150\,000 & $+45\,000$ & $-45\,000$ \\
9 -- Italia & 125\,000 & $+37\,500$ & $-37\,500$ \\
\bottomrule
\end{tabular}
\end{table}

\subsection{Capacidades de interconexión}

Los límites térmicos de las interconexiones se toman de los valores
públicos de \gls{NTC} comercial de \gls{ENTSOE} para 2024 y de las
especificaciones del proyecto CONECCMAR (Tabla~\ref{tab:lines}).

\begin{table}[h]
\centering
\caption{Capacidades térmicas de las interconexiones modelizadas.}
\label{tab:lines}
\begin{tabular}{lcl}
\toprule
Interconexión & Capacidad & Fuente \\
\midrule
España--Francia (CA) & \SI{2.8}{GW} & ENTSO-E 2024 \\
Francia--Italia (CA) & \SI{4.35}{GW} & ENTSO-E 2024 \\
España--Italia (\gls{HVDC}) & \SI{2.0}{GW} & Especificación CONECCMAR \\
Estaciones convertidoras & \SI{2.5}{GVA} & Sobredimensionada para soporte reactivo \\
\bottomrule
\end{tabular}
\end{table}

El límite térmico del cable \gls{HVDC} se programa con
$\SI{50}{MW}$ adicionales sobre los $\SI{2}{GW}$ nominales para absorber
las pérdidas $I^2R$ del cable, del orden de $\SI{30}{MW}$ a plena potencia.

\subsection{Costes marginales relativos}

Los coeficientes $c_{1,g}$ se eligen proporcionalmente al precio
mayorista medio horario de la electricidad en cada país en el periodo
2015--2026 (Open Power System Data). Normalizando España a $1{,}00$ se
obtienen los valores de la Tabla~\ref{tab:cost}. Se utiliza la media a
largo plazo en lugar de valores recientes porque el presente estudio
tiene carácter de planificación.

\begin{table}[h]
\centering
\caption{Coste marginal relativo utilizado por generador.}
\label{tab:cost}
\begin{tabular}{lcc}
\toprule
Nudo / país & Precio medio [EUR/MWh] & $c_1$ (relativo) \\
\midrule
1 -- España & 70{,}23 & 1{,}00 \\
7 -- Francia & 76{,}67 & 1{,}10 \\
9 -- Italia & 99{,}26 & 1{,}40 \\
\bottomrule
\end{tabular}
\end{table}

\section{Implementación}
\label{sec:impl}

El modelo se ha implementado en Python utilizando Gurobi~12 como solver
\gls{SOCP}. El programa carga los parámetros del caso de estudio,
construye la topología \gls{CFBM} introduciendo los nudos virtuales
asociados a cada convertidor, y formula el problema descrito en la
sección~\ref{sec:formulation}: variables $W$, flujos activos y reactivos
de rama, despacho de generadores y soporte reactivo auxiliar; balances
nodales; restricciones de cono rotado~\eqref{eq:soc}; modelo del
transformador no ideal; cotas de la fuente reactiva auxiliar~\eqref{eq:qz};
límites térmicos y de tensión.

Las tolerancias del solver se han ajustado para garantizar una brecha de
relajación numéricamente despreciable. En particular se utilizan
{\tt BarConvTol} = {\tt BarQCPConvTol} = $10^{-9}$,
{\tt FeasibilityTol} = {\tt OptimalityTol} = $10^{-7}$ y
{\tt NumericFocus} = 2. Una vez resuelto el problema, se recuperan las
tensiones complejas a partir de la matriz $W$ siguiendo el Algoritmo~1
de \cite{HigueraGutierrez2025}.

\section{Resultados y validación}
\label{sec:results}

Los resultados se contrastan frente al solver de flujo de cargas AC/DC
Newton--Raphson de \gls{MATPOWER}-\gls{FUBM}, motor numérico utilizado
por \gls{VFlexP}. La comparación no es directa: el solver Newton--Raphson
resuelve un flujo de cargas con consignas fijas de despacho activo en los
nudos PV (y el nudo slack absorbe el residuo), mientras que el \gls{OPF}
convexo optimiza el despacho dentro de los límites \eqref{eq:obj}. No
obstante, sirve como referencia sólida de validez física.

\subsection{Despacho}

La Tabla~\ref{tab:dispatch} compara el despacho por país. La desviación
en el agregado es del orden del $0{,}02\%$.

\begin{table}[h]
\centering
\caption{Despacho activo por país (MW).}
\label{tab:dispatch}
\begin{tabular}{lccc}
\toprule
País & NR (referencia) & CFBM--SOC & $\Delta$ \\
\midrule
España & 33\,766 & 33\,002 & $-2{,}3\%$ \\
Italia & 50\,000 & 49\,650 & $-0{,}7\%$ \\
Francia (slack) & 71\,325 & 72\,411 & $+1{,}5\%$ \\
\midrule
Total & 155\,091 & 155\,063 & $-0{,}02\%$ \\
\bottomrule
\end{tabular}
\end{table}

\subsection{Tensiones y ángulos}

Las tensiones coinciden con la referencia dentro de $\SI{0.003}{pu}$ en
todos los nudos físicos (Tabla~\ref{tab:voltages}). Los nudos 2 y 8 quedan
regulados a $\SI{1.0}{pu}$ por la acción de los convertidores. Los
ángulos siguen el mismo patrón cualitativo (España adelanta, Italia se
retrasa respecto al slack francés); las desviaciones absolutas se deben
a la libertad de redespacho del \gls{OPF} frente a las consignas fijas
del flujo de cargas.

\begin{table}[h]
\centering
\caption{Tensiones y ángulos: comparación con la referencia Newton--Raphson.}
\label{tab:voltages}
\begin{tabular}{cclcccc}
\toprule
Bus & País & Función & NR $|V|$ & CFBM $|V|$ & NR $\angle$ & CFBM $\angle$ \\
\midrule
1 & España & Generador & 1{,}0170 & 1{,}01642 & $+28{,}78^\circ$ & $+21{,}74^\circ$ \\
2 & España & Carga (regulado) & 1{,}0000 & 1{,}00000 & $+27{,}85^\circ$ & $+20{,}83^\circ$ \\
3 & --- & DC España & 1{,}0151 & 1{,}01790 & --- & $0$ \\
4 & --- & DC intermedio & 1{,}0150 & 1{,}01780 & --- & $0$ \\
5 & --- & DC Italia & 1{,}0000 & 1{,}00270 & --- & $0$ \\
6 & Francia & Nudo de cruce & 0{,}9927 & 0{,}99258 & $-0{,}40^\circ$ & $-0{,}41^\circ$ \\
7 & Francia & Generador (slack) & 1{,}0000 & 1{,}00000 & $0$ & $0$ \\
8 & Italia & Carga (regulado) & 0{,}9996 & 1{,}00000 & $-35{,}14^\circ$ & $-38{,}95^\circ$ \\
9 & Italia & Generador & 1{,}0123 & 1{,}01263 & $-34{,}46^\circ$ & $-38{,}27^\circ$ \\
\bottomrule
\end{tabular}
\end{table}

\subsection{Utilización de las interconexiones}

El enlace \gls{HVDC} España--Italia opera al $\SI{100}{\percent}$ de su
capacidad como restricción de desigualdad: el \gls{OPF} satura el cable
porque, dada la estructura de costes relativos (Italia es la zona más
cara), maximizar el flujo por el enlace es la opción económicamente
óptima. La interconexión Francia--Italia opera al $\SI{97}{\percent}$ y
la España--Francia al $\SI{87}{\percent}$ de su \gls{NTC}.

\subsection{Brecha de la relajación SOC}

En el óptimo, la brecha de cono $W_{ii}\,W_{jj} - |W_{ij}|^2$ es como
máximo $1{,}5 \times 10^{-7}$ en las parejas físicas, lo cual es
indistinguible del ruido numérico del solver. La relajación es por
tanto ajustada en este caso: no existe margen entre la solución
\gls{SOCP} y una solución factible de rango uno sobre la red física.
Las parejas asociadas a los transformadores no ideales del modelo
\gls{CFBM} exhiben una brecha de magnitud unitaria, pero ello forma
parte del diseño del modelo \cite{HigueraGutierrez2025}: en esas parejas
$W_{ij}$ no participa en el balance, y el paso de potencia se encauza a
través de variables $P$ y $Q$ específicas.

\subsection{Validación física}

Reconstruyendo las tensiones complejas a partir de la matriz $W$ por el
Algoritmo~1 de \cite{HigueraGutierrez2025}, los módulos recuperados
coinciden con $\sqrt{W_{ii}}$ con cinco decimales en todos los nudos
físicos. Esto confirma que la solución \gls{SOCP} es físicamente
realizable y de rango uno sobre la red.

\section{Conclusiones}
\label{sec:conclusions}

La relajación \gls{CFBM}-\gls{SOC} aplicada al caso CONECCMAR resuelve a
optimalidad en una sola pasada con Gurobi y es numéricamente ajustada,
con una brecha de cono máxima del orden de $10^{-7}$ en las parejas
físicas. Su despacho reproduce el del flujo de cargas Newton--Raphson de
\gls{MATPOWER}-\gls{FUBM} dentro del $0{,}02\%$ en el agregado, y las
tensiones recuperadas coinciden con la referencia dentro de
$\SI{0.003}{pu}$ en todos los nudos del sistema. El acuerdo es
suficientemente bueno como para considerar el método validado físicamente
para el caso estudiado.

Bajo el escenario de costes considerado, el \gls{OPF} satura el enlace
\gls{HVDC} España--Italia por motivos económicos y no por una restricción
impuesta a priori sobre la consigna del convertidor. La interconexión
Francia--Italia opera muy próxima a su \gls{NTC} comercial, lo que
confirma que la representación de los límites de transferencia
fronterizos es coherente con el comportamiento esperado del sistema en
un horizonte de planificación.

Por último, la solución \gls{SOCP} es de rango uno sobre la red física:
las tensiones complejas reconstruidas por el algoritmo de recuperación
coinciden con $\sqrt{W_{ii}}$ hasta cinco decimales en todos los nudos
físicos. Esto la habilita como punto inicial de calidad para los solvers
iterativos no convexos del flujo de cargas o del \gls{OPF} no lineal, lo
que abre la vía a su integración futura como herramienta interna de
\gls{VFlexP}, bien como solver convexo independiente, bien como
generador de un punto inicial cercano al óptimo global.

\bibliographystyle{IEEEtran}
\bibliography{BibSOCOPF}

\end{document}
